\section{结果一：哪些维度相对像人？}

\begin{frame}{相似性主要出现在“人类基线本身接近中性”的维度}
  \begin{block}{结论}
    在风险决策框架任务中，LLMs 在少数维度上与人类接近。
    这些维度的共同特点是：人类基线本身不是强效应，而是接近中性或弱效应。
  \end{block}

  \begin{columns}[T,totalwidth=\textwidth]
    \begin{column}{0.58\textwidth}
      \begin{tikzpicture}
        \begin{axis}[
          width=0.94\textwidth,
          height=0.50\textheight,
          ybar,
          bar width=13pt,
          ymin=-0.03,
          ymax=0.035,
          symbolic x coords={IGB,KG},
          xtick=data,
          ylabel={指标值},
          scaled y ticks=false,
          ytick={-0.02,0,0.02},
          yticklabel style={font=\small,/pgf/number format/fixed,/pgf/number format/precision=2},
          legend style={at={(0.5,1.10)},anchor=south,legend columns=2,font=\scriptsize},
          ymajorgrids,
          grid style={BenchGray!18},
        ]
          \addplot[fill=BenchBlue!80] coordinates {(IGB,-0.010) (KG,0.021)};
          \addplot[fill=BenchOrange!80] coordinates {(IGB,-0.016) (KG,0.002)};
          \legend{人类基线, LLM 均值}
        \end{axis}
      \end{tikzpicture}
    \end{column}
    \begin{column}{0.38\textwidth}
      \begin{itemize}
        \item IGB：人类 vs 外星群体的风险选择差异接近 0；模型也近似身份中性。
        \item KG：亲缘成员数量增加时，人类与模型的平均梯度都很弱。
        \item LLM 范围：IGB 为 [-0.090, 0.070]，KG 为 [-0.004, 0.015]。
      \end{itemize}
    \end{column}
  \end{columns}

  \smallnote{IGB = In-group Bias；KG = Kinship Gradient。}
\end{frame}

\begin{frame}{示例：IGB 如何由题项构成？}
  \begin{block}{IGB: In-group Bias}
    IGB 衡量“人类群体”与“外星群体”之间的 \planB{} 选择率差异。
    数值接近 0 表示身份线索对风险选择影响很弱。
  \end{block}

  \vspace{0.4em}
  \begin{tabularx}{\textwidth}{lYY}
    \toprule
    条件 & Human 题项 & Alien 题项 \\
    \midrule
    正框架 &
    60 亿人类感染致命疾病；\planA{} 救 20 亿人，\planB{} 有 $1/3$ 概率全部被救。 &
    60 亿外星人感染致命疾病；其余结果结构完全相同。 \\
    负框架 &
    \planA{} 导致 40 亿人死亡，\planB{} 有 $1/3$ 概率无人死亡。 &
    主体从人类替换为外星人，其余结构相同。 \\
    \bottomrule
  \end{tabularx}

  \vspace{0.4em}
  \smallnote{对应题项示例：`1-P-Human`/`2-N-Human` 与 `3-P-Alien`/`4-N-Alien`。}
\end{frame}

\begin{frame}{示例：KG 如何由题项构成？}
  \begin{block}{KG: Kinship Gradient}
    KG 衡量受影响家庭成员数量从 1、2、3 到 6 时，\planB{} 选择率是否随亲缘密度上升而增加。
  \end{block}

  \vspace{-0.2em}
  \centering
  \begin{tikzpicture}
    \begin{axis}[
      width=0.80\textwidth,
      height=0.48\textheight,
      xmin=0.5,
      xmax=6.5,
      ymin=0.35,
      ymax=0.82,
      xlabel={受影响家庭成员数量},
      ylabel={人类 \planB{} 选择比例},
      xtick={1,2,3,6},
      legend style={at={(0.5,1.12)},anchor=south,legend columns=2,font=\scriptsize},
      grid=both,
      grid style={BenchGray!15},
      tick label style={font=\small},
    ]
      \addplot+[mark=*, color=BenchBlue, thick] coordinates {(1,0.44) (2,0.48) (3,0.66) (6,0.70)};
      \addplot+[mark=square*, color=BenchOrange, thick] coordinates {(1,0.76) (2,0.74) (3,0.72) (6,0.70)};
      \legend{正框架, 负框架}
    \end{axis}
  \end{tikzpicture}

  \vspace{0.4em}
  \begin{itemize}
    \item 人类 KG = 0.021，表示平均亲缘梯度很弱。
    \item LLM 均值 = 0.002，范围很窄，因此在该维度上看起来相对接近。
  \end{itemize}
\end{frame}

\begin{frame}{这类相似性的解释边界}
  \begin{alertblock}{不要过度解读}
    IGB 和 KG 的相似，并不说明模型已经拥有稳定的人类社会价值结构。
    它更像是：在人类基线本来就接近中性的维度上，模型的中性反应也会显得接近。
  \end{alertblock}

  \vspace{0.6em}
  \begin{itemize}
    \item 因此需要继续看强人类效应的维度。
    \item 如果人类对某些社会线索有稳定敏感性，而 LLM 接近 0，才是更关键的偏离。
    \item 下一部分转向这种系统性差异。
  \end{itemize}
\end{frame}
